% !TEX program = xelatex
%----------------------------------------------------------------------------------------
%  NJUST Beamer Template
%  南京理工大学 Beamer 幻灯片模板
%  基于 SimplePlus Beamer Theme 修改
%  编译: XeLaTeX → Biber → XeLaTeX → XeLaTeX
%----------------------------------------------------------------------------------------

\documentclass[aspectratio=169,xcolor=dvipsnames]{beamer}

\usepackage{njustbeamer/beamerthemeSimplePlus}

\usepackage{tikz}
\usetikzlibrary{positioning}
\usepackage{fontspec}
\usepackage{xeCJK}
\usetikzlibrary{calc}

\usepackage{hyperref}
\usepackage{graphicx}
\usepackage{booktabs}
\usepackage{adjustbox}
\usepackage{pgfplots}

\usepackage[backend=biber,style=apa,sorting=none]{biblatex}
\addbibresource{njustbeamer/reference.bib}

%----------------------------------------------------------------------------------------
%  TITLE PAGE
%----------------------------------------------------------------------------------------

\title{南京理工大学 \\ 学术报告模板}
\subtitle{NJUST Academic Presentation Template}

\author{NLGer}

\institute
{
    XXX教研室 \\
    南京理工大学
}
\date{\today}

%----------------------------------------------------------------------------------------
%  PRESENTATION
%----------------------------------------------------------------------------------------

\begin{document}

\begin{frame}
    \titlepage
\end{frame}

{
    \setbeamertemplate{frametitle}{}

    \begin{frame}{目录}
    \begin{tikzpicture}[overlay, remember picture]
        \fill[Color3rd]
        (current page.north west) ++(0.20\paperwidth,-0.1\paperheight) rectangle ++(0.005\paperwidth,-0.8\paperheight);
        \node[
            rotate=90,
            anchor=center,
            text=Color1st,
            font=\bfseries\fontsize{36}{48}\selectfont
            ] at ([xshift=0.12\paperwidth]current page.west) {目录};
    \end{tikzpicture}

    \begingroup
    \leftskip=4cm
    \tableofcontents
    \endgroup
    \end{frame}
}

%------------------------------------------------
\section{引言}
%------------------------------------------------

\begin{frame}{项目符号}
    \begin{columns}[c]
        \column{.45\textwidth}
        这里可以放置左栏内容。该模板基于 SimplePlus 主题，适配南京理工大学视觉风格。
        \column{.45\textwidth}
        支持中文排版（需 XeLaTeX 编译），内置目录自动分节、参考文献引用等功能。
    \end{columns}
    \begin{tikzpicture}[overlay, remember picture]
        \coordinate (rectstart) at ([xshift=2em,yshift=1em]current page.south west);
        \fill[Color1st] (rectstart) rectangle ++(\textwidth-1em,30pt);
        \node[anchor=center,text=white,font=\large\itshape,text width=\textwidth,align=center
        ] at ($(rectstart)+(0.5*\textwidth,14pt)$) {南京理工大学 \quad XXX教研室};
    \end{tikzpicture}
\end{frame}

\begin{frame}{模块与列表}
    \begin{columns}[c]
        \column{.45\textwidth}
        \begin{block}{普通模块}
            \alert{强调文本} 这是一个普通模块。可用于展示重要概念或定义。
        \end{block}
        \begin{block}{项目列表}
            \begin{itemize}
                \item 第一个要点
                \item 第二个要点
                \item 第三个要点
            \end{itemize}
        \end{block}
        \column{.45\textwidth}
        \begin{block}{示例模块}
            这里展示示例模块的样式效果。
        \end{block}
    \end{columns}
\end{frame}

%------------------------------------------------
\section{核心内容}
%------------------------------------------------

\begin{frame}{公式与定理}
    \begin{theorem}[质能方程]
        $E = mc^2$
    \end{theorem}

    \begin{block}{要点}
        \begin{itemize}
            \item 定理环境自动编号并高亮
            \item 支持数学公式与交叉引用
        \end{itemize}
    \end{block}
\end{frame}

\begin{frame}{表格}
    \begin{table}
        \begin{tabular}{l l l}
            \toprule
            \textbf{方法} & \textbf{指标一} & \textbf{指标二} \\
            \midrule
            方法 A         & 0.003262         & 0.562           \\
            方法 B         & 0.015681         & 0.910           \\
            方法 C         & 0.009271         & 0.296           \\
            \bottomrule
        \end{tabular}
        \caption{结果对比}
    \end{table}
\end{frame}

\begin{frame}[fragile]{引用}
    使用 \verb|\cite| 命令引用文献：

    相关研究参见文献 \cite{p1}。
\end{frame}

%------------------------------------------------
\section{图片示例}
%------------------------------------------------

\begin{frame}{左文右图}
    \njtextimageright{%
        \textbf{文字说明}
        \begin{itemize}
            \item 使用 \texttt{\textbackslash njtextimageright} 命令
            \item 左侧放置文字内容
            \item 右侧放置配图
        \end{itemize}
    }{pic/01.jpg}{示例图片一}
\end{frame}

\begin{frame}{上文双图}
    \njtexttopimages{%
        \textbf{文字说明}
        \begin{itemize}
            \item 使用 \texttt{\textbackslash njtexttopimages} 命令
            \item 文字内容在上方
            \item 两张图片并排置于下方
        \end{itemize}
    }{pic/01.jpg}{示例图片一}{pic/02.jpg}{示例图片二}
\end{frame}

%------------------------------------------------
\section{总结}
%------------------------------------------------

\begin{frame}{多栏布局}
    \begin{columns}[c]
        \column{.45\textwidth}
        \textbf{标题}
        \begin{enumerate}
            \item 第一条
            \item 第二条
            \item 第三条
        \end{enumerate}
        \column{.45\textwidth}
        此处可放置总结文字或补充说明。
    \end{columns}
\end{frame}

\begin{frame}
    \Huge{\centerline{\textbf{谢谢！}}}
\end{frame}

% 参考文献页（取消注释以启用）
% \begin{frame}
%     \printbibliography
% \end{frame}

\end{document}
